\subsection{べき関数の広義積分}

前節では,\,$x=0$の近くで$\frac1{\sqrt{x}}$の積分は収束し,\,$\frac1x$の積分は発散した.どちらも限りなく大きくなる関数だが,大きくなる速さが違う.では,\,$\frac1{x^p}$では,\,$p$がいくつを境に収束と発散が分かれるのだろうか.原始関数が分かるので,直接計算して調べよう.後で他の関数の収束を判定するときは,このべき関数と比べることになる.

まず無限大での
\[
 \int_1^{\infty}\frac1{x^p}\dd x
\]
を調べる.\,$p\ne1$なら,まず有限の$R$まで積分して
\[
 \int_1^R x^{-p}\dd x=\frac{R^{1-p}-1}{1-p}.
\]
$p>1$なら$R^{1-p}\to0$なので収束し,値は$\frac1{p-1}$となる.\,$p<1$なら$R^{1-p}\to\infty$なので発散する.\,$p=1$では$\log R\to\infty$であり,やはり発散する.

これに対して,有限端点$0$の近くでは
\[
 \int_{\varepsilon}^{1}x^{-p}\dd x
 =\frac{1-\varepsilon^{1-p}}{1-p}\qquad(p\ne1)
\]
である.今度は$\varepsilon\to0+$なので,\,$p<1$なら$\varepsilon^{1-p}\to0$となって収束し,値は$\frac1{1-p}$になる.\,$p>1$なら$\varepsilon^{1-p}\to\infty$,\,$p=1$なら$-\log\varepsilon\to\infty$であり,いずれも発散する.

\begin{center}
\begin{tabular}{lll}
\toprule
積分 & 収束する条件 & 収束時の値\\
\midrule
$\int_1^{\infty}x^{-p}\dd x$ & $p>1$ & $\frac1{p-1}$\\[4pt]
$\int_0^1x^{-p}\dd x$ & $p<1$ & $\frac1{1-p}$\\
\bottomrule
\end{tabular}
\end{center}

同じべき関数でも,不等号の向きが逆になる.遠くで減る速さを調べる場合と,端点で大きくなる速さを調べる場合では,問題が異なるからである.どちらでも$\frac1x$が境目になり,そのときの原始関数$\log x$は$R\to\infty$でも$\varepsilon\to0+$でも絶対値が限りなく大きくなる.条件の向きは,\,$x^{1-p}$の指数$1-p$の符号から思い出せばよい.


